\chapter*{How to Work This Book}
\addcontentsline{toc}{chapter}{How to Work This Book}
\markboth{How to Work This Book}{How to Work This Book}

This book is organised by \emph{subject}: algebra, then functions, then
matrices, and so on, because that is the order in which the mathematics
builds on itself. A course is organised by \emph{examination topic}, and the
two orders are not the same. The second mathematics course of a business
administration or economics degree typically opens with multivariable
differentiation and expects everything in Parts~I to~IV to be already in
place; it then works through convexity, unconstrained optimization, implicit
functions, Lagrange, Kuhn--Tucker and linear programming, and closes with a
set of deliberately mixed exercises.

This chapter reconciles the two. It sets out ten stages in the order the
course actually needs them, says which chapters serve each stage, how much
work each is, and --- the part that matters most --- how to tell when a
stage is finished. Working the book front to back is perfectly reasonable if
you have the time. If you do not, the stages below are the efficient path.

\section*{First: find out where to start}

A course states its prerequisites but cannot enforce them, so the single
most common way to waste a semester is to begin at your syllabus's first
topic while missing something two volumes earlier. The following twelve
questions take about half an hour. Each one you cannot answer confidently,
\emph{in writing, without looking anything up}, sends you to a chapter.

\begin{enumerate}[leftmargin=2.0em,itemsep=0.35ex]
  \item Give the maximal domain of $\dfrac{\sqrt{x-1}}{\ln(x-2)}$.
        \hfill (Ch.~10)
  \item Sketch $y=2-\mathrm{e}^{-x}$ and give its asymptote.
        \hfill (Ch.~11, 12)
  \item Shade $\{(x,y): x^{2}+y^{2}\le 9,\ y\ge x\}$.
        \hfill (Ch.~13)
  \item Solve $x+2y-z=3$, $2x-y+z=5$, $3x+y+2z=14$ by elimination.
        \hfill (Ch.~19)
  \item Give the eigenvalues of
        $\left(\begin{smallmatrix}2&1\\1&2\end{smallmatrix}\right)$ and say
        whether it is positive definite. \hfill (Ch.~22, 23)
  \item Differentiate $x^{2}\mathrm{e}^{-x}$ and $\dfrac{\ln x}{x}$.
        \hfill (Ch.~29, 30)
  \item Evaluate $\displaystyle\lim_{x\to0}\frac{1-\cos x}{x^{2}}$.
        \hfill (Ch.~28, 35)
  \item Locate the extrema and the inflection of $x^{3}-3x^{2}$ and sketch
        it. \hfill (Ch.~36)
  \item Give $\dd{y}{x}$ on $x^{2}+xy+y^{2}=7$ at $(1,2)$.
        \hfill (Ch.~33)
  \item Sketch the level curves of $f(x,y)=x^{2}-y^{2}$ at
        $c=-1,0,1$. \hfill (Ch.~37)
  \item Give $\grad f$ for $f=x^{2}y^{3}$ at $(1,2)$ and
        $D_{\mathbf{u}}f$ towards $(3,4)$. \hfill (Ch.~38)
  \item Maximise $xy$ subject to $3x+4y=24$ and interpret the multiplier.
        \hfill (Ch.~57, 58)
\end{enumerate}

Questions 1--5 test the prerequisites, 6--9 the one-variable calculus,
10--12 the course proper. If you stumble in the first block, start at
Stage~1 and do not skip it; the later material will not stick on top of a
gap there.

\section*{How long this takes}

Worked honestly, a drill problem with a full solution takes four to six
minutes: read it, attempt it, check it, and note what went wrong. The ten
stages together are about $3\,200$ problems, which is \textbf{215 to 320
hours} --- a serious semester, and more than most people have.

So the stages below each carry a \textbf{core} and a \textbf{defer} list.
The core is what the syllabus actually leans on; taking only those chapters
brings the total to about $2\,480$ problems, or \textbf{165 to 250 hours}.
Two workable shapes:

\begin{itemize}[leftmargin=1.6em,itemsep=0.3ex]
  \item \textbf{A full semester}, 12 to 15 hours a week for fourteen weeks:
        Stages~1--4 in weeks 1--5, Stages~5--7 in weeks 6--8, Stages~8--9 in
        weeks 9--10, Stage~10 in weeks 11--13, and the mixed sets in
        week~14. This covers the core comfortably and most of the rest.
  \item \textbf{Four weeks left}, 20 to 30 hours a week: core chapters only,
        and within each chapter the \emph{Practice Problems} block only,
        skipping the Additional blocks entirely. Do every mastery test, and
        keep the mixed sets for the last four days: they are the only
        honest rehearsal.
\end{itemize}

One structural point that saves a lot of time: \textbf{Stage~3 can be
postponed}. Nothing in Stages~4 to~7 needs a matrix. Definiteness is first
required in Stage~8 and elimination in Stage~10, so if the calculus is the
urgent part, run Stages~1, 4, 5, 6, 7 first and fetch the matrices when
Stage~8 demands them.

\section*{The ten stages}

Each stage names the course topic it serves, the chapters that serve it, and
a \emph{mastery test} --- a specific thing you should be able to do before
moving on. The problem counts are the full drill; on a second pass, the
\emph{Additional Problems} block of each chapter alone is usually enough.

\subsection*{Stage 1: Functions, before anything else (Ch.~10--14, 400 problems)}

Syllabuses list this under prerequisites in a single sentence: you should
know what the graphs of the usual functions look like and be able to sketch
curves given by equations in two variables. That sentence hides a great deal.
Domain and range, the standard shapes with their poles and asymptotes,
shifting and stretching, inverses, and the regions that inequalities cut out
of the plane are used on nearly every page that follows.

\textbf{Mastery test.} Given $f(x)=\dfrac{x^{2}}{x-1}$, state the domain,
both asymptotes and the two extrema without differentiating anything you
have not already differentiated; and shade
$\{x^{2}+y^{2}\le16,\ y\ge0,\ x\le2\}$ and write the three inequalities back
out from your own picture.

\emph{Core:} Ch.~10--13. \emph{Defer:} Ch.~14 (trigonometry) --- the
prerequisites list it, but a business or economics examination almost never
uses it.
Do it only if the graphs of $\sin$, $\cos$ and $\tan$ are genuinely
unfamiliar.

\subsection*{Stage 2: The algebra underneath (Ch.~1--9, 275 problems)}

Only if Stage~1 was uncomfortable. This is manipulation: exponents,
logarithms, factoring, inequalities, systems, and the first-order-condition
algebra that every optimization problem eventually demands. Chapter~9 is the
one to do even if you skip the rest.

\textbf{Mastery test.} Isolate $K$ in $Y=AK^{\alpha}L^{1-\alpha}$, and solve
$\alpha y/(\beta x)=p_{x}/p_{y}$ together with $p_{x}x+p_{y}y=M$ without
hesitating.

\emph{Core:} Ch.~9, and whichever of Ch.~1--5 the diagnostic exposed.
\emph{Defer:} Ch.~7 (summation) and Ch.~8 (partial fractions) --- the first
matters for series, the second only for integration.

\subsection*{Stage 3: Matrices and linear systems (Ch.~15--24, 525 problems)}

Syllabuses name determinants, eigenvalues, definiteness, rank and linear
systems as review material. Two of these do real work later: definiteness is
the second-order condition of every multivariable optimum, and elimination
is how every linear system in the book is actually solved.

\textbf{Mastery test.} Row-reduce a $3\times3$ system to echelon form and
state which of the three outcomes it has; and classify
$\left(\begin{smallmatrix}2&-1\\-1&2\end{smallmatrix}\right)$ by both
eigenvalues and leading principal minors, getting the same answer twice.

\emph{Core:} Ch.~18, 19, 21, 22, 23 --- exactly the list a syllabus names:
determinants, elimination, rank, eigenvalues and definiteness.
\emph{Defer:} Ch.~15--17 unless matrix arithmetic is rusty, and Ch.~24
(applications), which is enrichment.

\subsection*{Stage 4: One-variable calculus, finished properly (Ch.~25--28 and 29--36, 706 problems)}

Sequences, series and limits first, then the differentiation rules,
L'H\^opital, and the complete curve discussion. The last of these,
\de{Kurvendiskussion}, is named explicitly in the prerequisites and is the
best single rehearsal of everything before it: domain, asymptotes,
monotonicity, curvature and a drawing, in one fixed order.

\textbf{Mastery test.} Run a full discussion of $f(x)=\dfrac{\ln x}{x}$ ---
domain, asymptotes, the maximum, the inflection --- and draw it; then
evaluate $\lim_{x\to0}(\cos x)^{1/x^{2}}$.

\emph{Core:} Ch.~28--31, 33, 35, 36. \emph{Defer:} Ch.~25--27 (sequences,
series, financial mathematics) unless your course examines them, Ch.~32
(Taylor) until Stage~8 needs the second-order intuition, and Ch.~34, which
overlaps with Stage~8.

\subsection*{Stage 5: The plane in two variables (Ch.~37, 92 problems)}

Syllabus topic: domain, level curves and range for functions of two
variables (\de{Definitionsmenge, Niveaulinien, Wertemenge}). Surfaces,
contour maps, sub- and superlevel sets. This is short and is where most of the geometry of the rest of the
book is decided, so it deserves careful work.

\textbf{Mastery test.} Match a contour map to its surface and justify it by
the \emph{spacing} of the contours; and read the critical points of
$x^{3}-3x+y^{2}$ off its contour map, classifying each, before computing
anything.

\subsection*{Stage 6: Multivariable differentiation (Ch.~38--41, 340 problems)}

Syllabus topics: multivariable differentiation and the chain rule
(\de{Differentialrechnung in mehreren Variablen}, \de{Kettenregel}).
Partial derivatives, the gradient,
directional derivatives, the total differential, the multivariable chain
rule, second-order partials and the Hessian.

\textbf{Mastery test.} State why $\grad f$ is perpendicular to a level
curve, in two sentences, without looking it up --- everything in Stages~8
and~9 leans on it.

\emph{Core:} all four. This is the centre of the course and nothing here is
optional.

\subsection*{Stage 7: Convexity (Ch.~52--53, 113 problems)}

Syllabus topic: convexity (\de{Konvexit\"at}). Convex sets, convex and
concave functions, and quasiconcavity. Short, and it pays for itself twice, since
convexity is what makes a first-order condition sufficient later.

\textbf{Mastery test.} Decide convexity of a function from its Hessian and
of a set from the line-segment test, and explain why $xy$ is quasiconcave on
the positive quadrant but not concave.

\subsection*{Stage 8: Unconstrained optimization and the envelope (Ch.~54--56, 111 problems)}

Syllabus topic: unconstrained optimization in several variables
(\de{Optimierung ohne Nebenbedingungen}). Critical points, the Hessian
test, comparative statics and the envelope theorem.

\textbf{Mastery test.} Find and classify every critical point of
$x^{3}-3xy+y^{3}$; then estimate the change in an optimal value when a
parameter shifts, \emph{without} re-solving the problem.

\emph{Core:} all three. Ch.~56 (comparative statics and the envelope) is the
one students skip and then meet again, unrecognised, inside Stage~10.

\subsection*{Stage 9: Implicit and inverse functions (Ch.~33 and 42, 198 problems)}

Syllabus topic: inverse and implicit functions
(\de{Inverse und implizite Funktionen}). Chapter~33 is the mechanics in one
variable, Chapter~42 the two named theorems --- implicit and inverse ---
together with the Jacobian, homogeneity and Euler's theorem. The pair are
one statement twice: the implicit theorem asks whether an equation can be
solved for a variable, the inverse theorem whether a map can be undone, and
both turn on a derivative (or a determinant) being non-zero. The thing to
take away is where they \emph{fail}: at a vertical tangent, a
self-crossing, or a vanishing Jacobian determinant.

\textbf{Mastery test.} Given the picture of an implicit curve, name every
point where $y$ cannot be written as a function of $x$, justify it from the
picture, and confirm it by computing $F_{y}=0$ together with $F=0$.

\emph{Core:} both. Within Ch.~42, the two theorems and the Jacobian are the
chapter and cannot be skipped; homogeneity and Euler's theorem can wait
until you meet constant returns to scale.

\subsection*{Stage 10: Constrained optimization and linear programming (Ch.~57--60 and 61--68, 470 problems)}

The final stage, and the largest block of the examination: Lagrange with its
first- and second-order conditions and the interpretation of the multiplier;
the bordered Hessian; Kuhn--Tucker with complementary slackness; and then
linear programming end to end. On most syllabuses these two topics carry
more examination weight than any other pair, so budget accordingly.

\textbf{Mastery test.} Solve a constrained problem by Lagrange, certify it
with the bordered Hessian, and re-derive it by substitution; then solve a
Kuhn--Tucker problem by systematic case analysis, and read an optimal
simplex tableau for the shadow prices.

\emph{Core:} all of it. If time forces a choice, Ch.~57--60 outrank the
linear-programming chapters: Lagrange and Kuhn--Tucker are the techniques a
business or economics examination reaches for most often, and they carry
into later courses in a way the simplex tableau does not.

\section*{Finally: the mixed sets}

Everything above drills one technique at a time, which is the right way to
learn a method and the wrong way to sit an examination --- a chapter heading
tells you what to do before you have read the question. The first three
sets of Part~IX (Ch.~69--71) announce nothing. Work one in a single sitting under a
time limit, then for every problem you got wrong write down not the
arithmetic but the \emph{decision}: which signal in the wording should have
pointed at which method? That note is worth more than redoing the
computation, and it is the thing a syllabus's closing set of cross-topic
problems (\emph{Kapitel\"ubergreifende Aufgaben}) exists to train.
The fourth set (Ch.~72) trains the other half of that skill: it takes one
function of two variables through the whole chain of questions, from domain
and range to gradient, Hessian, contour sets and a checklist of properties,
so that each answer has to agree with the ones before it.

\section*{Mapping a syllabus to this book}

Business and economics mathematics courses usually hand out an exercise
collection organised by topic, and the topics are remarkably standard from
one degree to the next. Below are the ones such a collection almost always
contains, with the German headings under which they are usually taught in
the German-speaking countries. Match your own syllabus against it: the
wording will differ, the topics will not.

\begin{center}\footnotesize
\begin{tabular}{@{}lllr@{}}
\textbf{Syllabus topic} & \textbf{German heading} &
  \textbf{Chapters} & \textbf{Problems}\\
\hline
One-variable functions, revision & \de{Wiederholung eindimensionale Funktionen}
  & 10--14, 29--36 & 986\\
Inequalities in two variables & \de{Ungleichungen in zwei Variablen}
  & 13 & 97\\
Domain, level curves, range & \de{Definitionsmenge, Niveaulinien, Wertemenge}
  & 10, 37 & 162\\
Multivariable differentiation & \de{Differentialrechnung in mehreren Variablen}
  & 38, 39, 41 & 278\\
Chain rule & \de{Kettenregel} & 40 & 62\\
Convexity and the Hessian & \de{Konvexit\"at und Hessematrix}
  & 39, 52, 53 & 198\\
Unconstrained optimization & \de{Optimierung ohne Nebenbedingungen}
  & 54, 55, 56 & 111\\
Inverse and implicit functions & \de{Inverse und implizite Funktionen}
  & 33, 42 & 198\\
Equality constraints (Lagrange) & \de{Optimierung mit Gleichungen}
  & 57, 58, 59 & 184\\
Nonlinear programming (KKT) & \de{Nichtlineare Programmierung}
  & 60 & 78\\
Linear programming & \de{Lineare Programmierung} & 61--68 & 208\\
Cross-topic problems & \de{Kapitel\"ubergreifende Aufgaben} & 69--72 & 65\\
\end{tabular}
\end{center}

Two chapters sit outside that table and are worth knowing about. Part~VI
(integration, Ch.~43--50) is not on a typical second-semester syllabus at
all, but it is on most business and economics degrees somewhere, and consumer
surplus and present value live there. Chapter~27, financial mathematics, is the same
case.

\section*{What this book does not do}

These limits are worth knowing before you rely on the book.

\begin{itemize}[leftmargin=1.6em,itemsep=0.4ex]
  \item \textbf{No $\varepsilon$--$\delta$.} Limits are established with the
        limit rules, L'H\^opital and algebra, never with
        $\varepsilon$-arguments. That matches a mathematics-for-economists
        syllabus, which asks you to \emph{compute} limits rather than prove
        them exist. If your course does demand the rigorous definition, it
        is the one thing you will have to get elsewhere.
  \item \textbf{Proofs are demonstrations, not verifications.} Results are
        derived far enough that they are believable and usable --- the
        gradient's perpendicularity, $AC'=(MC-AC)/q$, Euler's theorem --- but
        this is not an analysis text and does not pretend to be.
  \item \textbf{Integration is a bonus, not a requirement.} Part~VI and
        Ch.~27 sit outside a typical second-semester syllabus. They are here
        because a business or economics degree wants them somewhere. Skip
        them unless your course lists them.
  \item \textbf{No statistics, no econometrics, no differential equations.}
        Ch.~24 touches OLS only as a matrix identity.
\end{itemize}

\section*{Practical notes}

\begin{itemize}[leftmargin=1.6em,itemsep=0.4ex]
  \item \textbf{Start each chapter with its worked example.} Every chapter
        carries one, between the rule boxes and the problems: a single
        representative problem solved end to end, with the traps named where
        they occur and every answer checked. Read it with a pen, reproducing
        each step before looking at it, and only then start the drill. It is
        the shortest route from ``I have read the rules'' to ``I can use
        them'', and the \emph{Trap} notes in it are the errors that actually
        cost marks.
  \item \textbf{Work the problem before reading the solution.} Every
        solution here is complete, which makes it very easy to read one and
        feel that you could have produced it. You could not, until you have.
  \item \textbf{When you are stuck, read only the first line.} The solutions
        in Part~IX open by naming the technique and nothing else; elsewhere
        the first sentence usually does the same. Read that one line, close
        the book, and try again. Most failures are failures of
        \emph{selection}, not of execution, and reading the whole solution
        hides which of the two yours was.
  \item \textbf{Draw everything.} The course asks you to read properties off
        pictures --- convexity from a surface, critical points from a
        contour map, binding constraints from a shaded region. That skill is
        only built by drawing, not by reading drawings.
  \item \textbf{Use the back matter.} The \emph{Formula Summary} is what to
        revise from in the last week, and the \emph{German--English
        Glossary} is there for degrees that teach from English sources and
        examine in German.
  \item \textbf{When stuck, go one stage back.} Almost every difficulty in
        Stages~6 to~10 is really a gap in Stage~1, 3 or 4 --- an
        unrecognised domain, a definiteness test misapplied, a derivative
        taken wrongly.
\end{itemize}
